\documentclass[12pt]{article}
\usepackage[T1]{fontenc}
\usepackage[utf8]{inputenc}
\usepackage{lmodern}
\usepackage{textcomp}
\usepackage{enumitem}
\usepackage{geometry}
\geometry{margin=1in}
\begin{document}
\begin{center}
Answer Key - Health Sciences - Graduate Level Assessment 16 \
\small (Answers are ordered exactly as in the question paper.)
\end{center}

\section*{Multiple Choice}
\begin{enumerate}[label=\arabic*.]
\item \textbf{Answer:} A
\\ \textit{Options:} A) ensure the emergency team/services are called.; B) administer an electric shock using a defibrillator.; C) give 30 chest compressions.; D) check the notes to see if the patient has a DNAR order.; E) rush the patient to the emergency room.
\\ \textit{Note:} The correct answer is to ensure the emergency team/services are called because rapid activation of emergency medical services is critical in cardiac arrest situations to facilitate timely advanced care and improve the chances of survival.
\item \textbf{Answer:} A
\\ \textit{Options:} A) Palmitic acid; B) Pyruvate; C) Galactose; D) Glycerol; E) Propionic acid
\\ \textit{Note:} Palmitic acid cannot be a substrate for gluconeogenesis because it is a fatty acid that undergoes beta-oxidation to produce acetyl-CoA, which cannot be converted back to glucose.
\item \textbf{Answer:} C
\\ \textit{Options:} A) Superiority; B) Guilt; C) Personal control; D) Material wealth; E) Humor
\\ \textit{Note:} Individuals with a sense of personal control tend to engage more deeply with their life narratives, leading to a greater likelihood of constructing coherent life stories that reflect their experiences and agency in shaping their lives.
\item \textbf{Answer:} A
\\ \textit{Options:} A) Starting a new business; B) Working part-time; C) Mountain climbing; D) Household chores; E) Travel
\\ \textit{Note:} Starting a new business is not commonly cited as a frequent activity among retirees, who typically engage in leisure activities, socializing, or volunteering rather than entrepreneurship.
\item \textbf{Answer:} B
\\ \textit{Options:} A) 0.001; B) 0.01; C) 0.1; D) 0.15; E) 0.0011
\\ \textit{Note:} The expected frequency of affected females is 0.01 because X-linked recessive conditions require both X chromosomes to carry the recessive allele for females to be affected, so with a male frequency of 0.10, the carrier frequency in females is 0.30, leading to 0.01 affected females (0.30 * 0.10).
\end{enumerate}

\section*{True / False}
\begin{enumerate}[label=\arabic*.]
\setcounter{enumi}{5}
\item \textbf{Answer:} True
\\ \textit{Options:} A) True; B) False
\\ \textit{Note:} Vitamin D primarily functions in calcium metabolism and bone health rather than serving as a lipid-soluble antioxidant, which is a role typically associated with other compounds like vitamin E.
\item \textbf{Answer:} True
\\ \textit{Options:} A) True; B) False
\\ \textit{Note:} Macrostomia is indeed the result of the failure of fusion of the left and right maxillary processes during embryonic development, leading to an abnormal widening of the oral cavity.
\item \textbf{Answer:} False
\\ \textit{Options:} A) True; B) False
\\ \textit{Note:} Cerebrospinal fluid circulates within the subarachnoid space, which is located between the arachnoid mater and the pia mater, not between the skull and the dura mater.
\item \textbf{Answer:} True
\\ \textit{Options:} A) True; B) False
\\ \textit{Note:} Inadequate MHC Class I presentation reduces the ability of CD 8 T cells to recognize and respond to HIV-infected cells, leading to impaired immune responses and persistence of the virus in chronically infected individuals.
\item \textbf{Answer:} True
\\ \textit{Options:} A) True; B) False
\\ \textit{Note:} Damage to the third cranial nerve, which innervates most of the extraocular muscles and controls pupil constriction, leads to impaired muscle function resulting in pupillary dilatation (due to loss of sphincter muscle control) and downward strabismus (as the unopposed action of the inferior oblique and superior oblique muscles causes the eye to deviate downwards).
\end{enumerate}

\section*{Long Form Response}
\begin{enumerate}[label=\arabic*.]
\setcounter{enumi}{10}
\item \textbf{Answer:} Genetic factors play a significant role in longevity, influencing physiological processes and susceptibility to age-related diseases. Individuals with "good genes" may possess advantageous alleles that enhance cellular repair mechanisms, reduce inflammation, and promote overall health. However, these genetic predispositions interact with lifestyle factors such as mental stimulation, exercise, and diet, which are crucial for achieving a long and healthy life. Mental stimulation fosters cognitive resilience, exercise promotes cardiovascular health and metabolic efficiency, and a balanced diet provides essential nutrients that support cellular function. Collectively, these elements create a synergistic effect, where favorable genetics can be optimized through healthy lifestyle choices, ultimately contributing to an increased lifespan and improved quality of life.
\\ \textit{Note:} correct because it effectively highlights the interplay between genetic factors and lifestyle choices-such as mental stimulation, exercise, and diet-in determining longevity, emphasizing that while genetics provide a foundation for health, lifestyle significantly modulates these genetic influences.
\item \textbf{Answer:} Alcoholic myopathy is a common condition among individuals with alcohol use disorder, affecting approximately 40 to 60\% of this population. Contributing factors include nutritional deficiencies, particularly of thiamine and other essential nutrients, as well as the toxic effects of alcohol on muscle tissue and metabolism. These factors lead to muscle weakness and atrophy, significantly impacting physical health and recovery efforts. Consequently, treatment approaches must address both the cessation of alcohol consumption and the nutritional rehabilitation of affected individuals to enhance recovery outcomes. Understanding the prevalence and implications of alcoholic myopathy is essential for developing effective interventions in clinical settings.
\\ \textit{Note:} correct because it accurately identifies the high prevalence of alcoholic myopathy among individuals with alcohol use disorder and discusses the significant contributing factors, such as nutritional deficiencies and the toxic effects of alcohol, which are critical for understanding the implications for effective treatment and recovery.
\end{enumerate}

\vfill
\noindent\textit{End of Answer Key}
\end{document}